% output=pdftex

\environment mstyle

\definefont [BigFont] [RegularBold at 75pt]
\definecolor[BigColor][s=.9]

\startuseMPgraphic{page}
  StartPage ;

  path p[] ; color c[] ; picture pic ;

  % yellow instead of red also looks ok (but does not match the style)

  c[1] := \MPcolor{blue} ;
  c[2] := \MPcolor{red} ;
  c[3] := .1white ;

  p[1] := (0,0) -- (0,100) ;
  p[2] := (0,50) {right} .. (50,75) .. {left}  (0,100) ;
  p[3] := (10,0) {right} .. {left} (30,25) .. {right} (50,50) ;
  p[4] := (25,50) -- (25,0) ;
  p[5] := (25,0) {left} .. (10,12.5) .. {right} (25,25) ;
  p[6] := (35,0) -- (35,35) {up} .. {right} (50,50) ;
  p[7] := (35,25) -- (50,25) ;

  pic := image
    ( for i := 1 upto 100 :
        for j := 1 upto 2 :
          draw p[j] randomized 15
            withcolor c[if odd i: 1 else: 2 fi] randomized (.5,.8) ;
        endfor ;
        for j := 3 upto 3 :
          draw p[j] randomized 10 withcolor c[1] randomized (.5,.8) ;
        endfor ;
        for j := 4 upto 5 :
          draw p[j] randomized 10 withcolor c[2] randomized (.5,.8) ;
        endfor ;
        for j := 6 upto 7 :
          draw p[j] randomized 10 withcolor c[2] randomized (.5,.8) ;
        endfor ;
      endfor ; ) ;

  pic := pic shifted - llcorner pic ;
  pic := pic xysized (bbwidth(Page)-20pt,bbheight(Page)-20pt) ;
  pic := pic shifted (10pt,10pt) ;

  fill Page withcolor c[3] ;
  draw pic ;

  StopPage ;
\stopuseMPgraphic

\starttext

\setuplayout[page]

\definelayer[page][width=\paperwidth,height=\paperheight]

\startstandardmakeup
  \setlayer[page]
    {\useMPgraphic{page}}
  \setlayerframed[page]
    [preset=rightbottom]
    [width=fit,height=fit,
     align=left,offset=10pt,
     frame=off,foregroundcolor=BigColor]
    {\BigFont \setstrut
     \strut Weaving
     \vskip-15pt
     \strut PS into PDF}
  \flushlayer[page]
\stopstandardmakeup

\setuplayout[reset]

\section {Introduction}

This manual will describe the \PSTOPDF\ conversion tool,
written in \RUBY\ on top of \GHOSTSCRIPT. When going from
\POSTSCRIPT\ to \PDF, you can use \DISTILLER\ and a few
other commercial applications, as well as the free
\GHOSTSCRIPT\ tool. Each tool has its advantages and in
professional workflows they will be used intermixed.

The \PSTOPDF\ converter evolved out of the converter built
into \TEXUTIL, a file utility that comes with \CONTEXT.
That converter itself was inspired by the \PDF\ to \PS\
pipeline posted by Sebastian Rahtz in the early days of
\PDFTEX.

Where \TEXUTIL\ goes beyond piping by applying some filters
to the \POSTSCRIPT\ code, \PSTOPDF\ adds even more
filtering and cleanup of \PS\ code, especially application
specific directives that fall outside the scope of
standards and converters.

The extensions were driven by the fact that in out projects
we were confronted with rather messy \POSTSCRIPT\ and
cleanup was needed anyhow. Not only the content can be a
problem, also filenames and suffixes can be messed up. This
is why \PSTOPDF\ is accompanied by a series of other image
handling tools, described elsewhere.

\section {The command line}

Because this tool is part of the \EXAMPLE\ toolkit that
comes with \CONTEXT, it's control may look strange but suits
this framework.

\starttabulate[|Tl|p|]
\NC --debug      \NC shows a lot of extra info                 \NC \NR
\NC --nopipe     \NC avoids piping, handy for buggy platsforms \NC \NR
\NC --method     \NC 1--5, more on that later                  \NC \NR
\NC --offset     \NC an additional offset to cropped images    \NC \NR
\NC --prefix     \NC tokens to prepend to teh output filename  \NC \NR
\NC --suffix     \NC tokens to append to teh output filename   \NC \NR
\NC --inputpath  \NC the path where input files are looked up  \NC \NR
\NC --outputpath \NC the path where output files are placed    \NC \NR
\stoptabulate

An example of usage is:

\starttyping
pstopdf --method=2 *.eps *.ai
\stoptyping

\section {Watched folders}

You can set up a watched folder (aka hot folder) structure.
This has the advantage that you only have to install one
instance of \PSTOPDF\ in a network.

\starttabulate[|Tl|p|]
\NC --watch      \NC folder to watch for files to convert     \NC \NR
\NC --delay      \NC interval between watch scans             \NC \NR
\stoptabulate

Say that we called \PSTOPDF\ as follows:

\starttyping
pdtopdf --watch=e:/pstopdf/watch --delay=2
\stoptyping

The program will look in subfolders of the \type {watch}
path for user subpaths, like:

\starttyping
.../pstopdf/watch/hagen
\stoptyping

Results will be placed in:

\starttyping
.../pstopdf/watch/hagen/result
\stoptyping

A converted file ends up in:

\starttyping
e:/pstopdf/watch/hagen/done
\stoptyping

You can control the conversion by putting the files into the
automatically created substructure:

\starttyping
.../pstopdf/watch/hagen/raw
.../pstopdf/watch/hagen/bound
.../pstopdf/watch/hagen/crop
\stoptyping

There is an bonus subpath:

\starttyping
.../pstopdf/watch/hagen/bitmap
\stoptyping

When bitmap images are placed on this path, they are
converted to \PDF\ using \IMAGEMAGICK. The compression and
quality chosen roughly depends on the size of the file.

\section {Batch control}

{\em To do: Here we will discuss how you can set up
conversion presets and or use \EXAMPLEX\ to control
\PSTOPDF.}

\section {Configuration}

There is not much configuration needed. First of all you
need to install \RUBY, \GHOSTSCRIPT, and optionally
\IMAGEMAGICK. Then you must locate \PSTOPDF\ in your \TEX\
tree. You can either call this script directly or call it
by passing the filename to \RUBY. This depends on your
operating system and configuration.

\starttyping
.../tex/texmf-local/context/ruby
\stoptyping

Details of the environment can be set up in the \EXAMPLE\
configuration files, located in (for instance):

\starttyping
.../tex/texmf-local/context/examplap/scripts
\stoptyping

The main configuration file is called \type {example.exa},
which calls up a few others. Don't mess around with the file
\type {basic.exa}, but stick to adapting \type {path.exe} or
\type {local.exa} if needed.

\section {Methods}

Currently there three conversion methods for \EPS\ images,
and two for \PDF\ images:

\starttabulate[|l|lT|lT|p|]
\NC \bf method \NC \bf suffix \NC prefix \NC action \NC \NR
\NC 1 \NC eps \NC         \NC leave boundingbox as it is            \NC \NR
\NC 2 \NC eps \NC         \NC crop image to internal bounding box   \NC \NR
\NC 3 \NC eps \NC         \NC crop image to analyzed bounding box   \NC \NR
\NC 4 \NC pdf \NC lowres- \NC downsample image to screen resolution \NC \NR
\NC 5 \NC pdf \NC normal- \NC cleanup image to prepress defaults    \NC \NR
\stoptabulate

When no output filename is specified, the prefix is appended to
the filename.

\section {The GUI}

We've written an interface to this tool, called \type
{pdtopdf.pdf}. You can find the interface in the path:

\starttyping
.../tex/texmf-local/context/examplap/gui
\stoptyping

The interface is more or less self documenting.

\stoptext